\chapter{Vollständige Legende der Simulation}

\section{Einleitung}

Dieser Anhang enthält die vollständige Ausgabelegende der IWT-Simulation.

Die Legende beschreibt jede Zeile der Simulationsausgabe und ihre physikalische Bedeutung.

Die Simulation ist der Beweis, dass die IWT funktioniert.

\section{Die Ausgabe der Simulation}

Die folgende Legende beschreibt die vollständige Ausgabe der Simulation.

\subsection{Zeile 1: Kernstatistik}

\begin{verbatim}
#: 499                       - Aktuelle Iteration (Zeitschritt)
|Q_max|: 240806.394          - Maximalwert des Bohm-Potentials Q
                                (je grösser, desto stärker die Strukturbildung)
I_total: 2.955               - Summe der Amplituden |I| über alle Knoten
                                (Mass für die gesamte Informationsmenge)
I_min: 0.000                 - Minimale Amplitude |I| (Vakuum)
I_max: 0.063                 - Maximale Amplitude |I| (Struktur)
Deviation: 6.214             - Abweichung von der Anfangsbedingung
                                (sollte stabil sein)
\end{verbatim}

\subsection{Zeile 2: Energiestatistik}

\begin{verbatim}
Summe |I|^2: 2.955           - Summe der Betragsquadrate |I|^2
                                (MUSS konstant sein = Axiom 2)
Info_Dev: 6.214              - Relative Abweichung von Summe |I|^2
                                (sollte nahe 0 sein)
\end{verbatim}

\subsection{Zeile 3: Fluktuationsstatistik}

\begin{verbatim}
xi_real_mean: -0.003         - Mittelwert der Realteil-Fluktuationen
xi_imag_mean: -0.003         - Mittelwert der Imaginärteil-Fluktuationen
xi_real_sigma: 0.004         - Standardabweichung der Realteil-Fluktuationen
xi_imag_sigma: 0.004         - Standardabweichung der Imaginärteil-Fluktuationen
scale: 0.707                 - Skalierungsfaktor der Fluktuationen
                                (aus der Theorie: sqrt(hbar/(2*T)))
\end{verbatim}

\subsection{Zeile 4: Summe |I|\texttwosuperior Historie}

\begin{verbatim}
Summe |I|^2 (letzte 10): 3.0 ... - Die letzten 10 Werte von Summe |I|^2
                                (sollten konstant sein = Axiom 2)
\end{verbatim}

\subsection{Zeile 5: I\_max Historie}

\begin{verbatim}
I_max (letzte 10): 0.1 ...   - Die letzten 10 Werte von I_max
                                (sollten stabil sein)
\end{verbatim}

\subsection{Zeile 6: Realteil-Statistik}

\begin{verbatim}
Re_mean: 0.022               - Mittelwert von Re(I) über alle Knoten
Re_min: 0.000                - Minimaler Realteil (Vakuum)
Re_max: 0.063                - Maximaler Realteil (Struktur)
                                (je grösser die Differenz Re_max - Re_min,
                                 desto stärker die Struktur)
\end{verbatim}

\subsection{Zeile 7: Imaginärteil-Statistik}

\begin{verbatim}
Im_mean: -0.000              - Mittelwert von Im(I) über alle Knoten
Im_min: -0.000               - Minimaler Imaginärteil
Im_max: 0.000                - Maximaler Imaginärteil
                                (sollte symmetrisch um 0 sein)
\end{verbatim}

\subsection{Zeile 8: Phasen-Statistik}

\begin{verbatim}
phi_mean: -0.000             - Mittelwert der Phase (in Grad)
phi_min: -0.005              - Minimale Phase
phi_max: 0.000               - Maximale Phase
                                (sollte symmetrisch um 0 sein)
\end{verbatim}

\subsection{Zeile 9: Fluss-Statistik}

\begin{verbatim}
J_mean: -0.000               - Mittelwert des Informationsflusses
J_min: -0.001                - Minimaler Fluss
J_max: 0.003                 - Maximaler Fluss
                                (Mittelwert sollte approx 0 sein = Erhaltung)
\end{verbatim}

\subsection{Zeile 10: Bohm-Potential Statistik}

\begin{verbatim}
Q_mean: -2022.872            - Mittelwert des Bohm-Potentials Q
                                (negativ = bindend / strukturierend)
\end{verbatim}

\subsection{Zeile 11: Realteil-Fluktuationen}

\begin{verbatim}
xi_r_mean: -0.003            - Mittelwert der Realteil-Fluktuationen
xi_r_min: -0.017             - Minimaler Wert der Realteil-Fluktuationen
xi_r_max: 0.001              - Maximaler Wert der Realteil-Fluktuationen
\end{verbatim}

\subsection{Zeile 12: Imaginärteil-Fluktuationen}

\begin{verbatim}
xi_i_mean: -0.003            - Mittelwert der Imaginärteil-Fluktuationen
xi_i_min: -0.017             - Minimaler Wert der Imaginärteil-Fluktuationen
xi_i_max: 0.001              - Maximaler Wert der Imaginärteil-Fluktuationen
\end{verbatim}

\section{Die physikalische Interpretation}

Die Werte der Simulation haben eine klare physikalische Bedeutung:

\begin{enumerate}
\item \textbf{Summe |I|\texttwosuperior{} ist konstant (2.955):}
      \begin{itemize}
      \item Axiom 2 (Informationserhaltung) ist erfüllt.
      \item Die Simulation ist stabil.
      \end{itemize}

\item \textbf{I\_max > I\_min (0.063 > 0.000):}
      \begin{itemize}
      \item Es gibt Struktur im System.
      \item Die Struktur ist stabil.
      \end{itemize}

\item \textbf{Q\_max ist gross (240806):}
      \begin{itemize}
      \item Das Bohm-Potential ist aktiv.
      \item Es organisiert die Struktur.
      \end{itemize}

\item \textbf{Q\_mean ist negativ (-2022.872):}
      \begin{itemize}
      \item Das Bohm-Potential ist bindend.
      \item Es hält die Struktur zusammen.
      \item Das ist Gravitation.
      \end{itemize}

\item \textbf{J\_mean approx 0:}
      \begin{itemize}
      \item Der Fluss ist im Gleichgewicht.
      \item Es gibt keinen Netto-Transport von Information.
      \end{itemize}

\item \textbf{Fluktuationen sind symmetrisch (xi\_r\_mean approx xi\_i\_mean):}
      \begin{itemize}
      \item Die Fluktuationen sind isotrop.
      \item Es gibt keine Vorzugsrichtung.
      \end{itemize}
\end{enumerate}

\section{Das Fazit der Simulation}

Die Simulation läuft stabil.

Axiom 2 ist erfüllt.

Strukturen existieren.

Das System ist im Gleichgewicht.

Alle Werte sind konsistent.

Die Simulation ist ein physischer Beweis, dass die IWT funktioniert.

\section{Zusammenfassung}

\begin{itemize}
\item Die vollständige Ausgabe der Simulation ist dokumentiert.
\item Jede Zeile hat eine klare physikalische Bedeutung.
\item Die Werte bestätigen die Theorie.
\item Die Simulation ist der Beweis, dass die IWT funktioniert.
\end{itemize}